\documentclass[11pt,a4paper]{article}
\usepackage[utf8]{inputenc}
\usepackage[T1]{fontenc}
\usepackage[spanish]{babel}
\usepackage[margin=2cm]{geometry}
\usepackage{booktabs,hyperref,array}
\hypersetup{colorlinks=true,urlcolor=blue}
\setlength{\parskip}{5pt}\setlength{\parindent}{0pt}
\begin{document}
\begin{center}{\Large Escalamiento interno y valor del generador}\\[4pt]
Exploración independiente para OMS --- octubre de 2026\end{center}
\textbf{Decisión: mantener OMS v1.3.5 sin cambios.} El escalamiento externo
produce beneficios en algunos casos con parámetros normales, pero no aparece
una ventaja atribuible a información adicional ni una regla predictiva validada.

\textbf{Diseño.} Seis modelos SG--Ter--Mer de 360 períodos, 6120 variables,
10080 filas y 1440 enteras cada uno. Descubrimiento: 001,014,042, elegidos por
resultados históricos; reserva: 002,003,004, fijados antes de optimizar. Se
conserva la integralidad original. No se usan los tres MPS continuos de Benders.
Se exportan MPS nativos desde el ejecutable existente, sin modificar fuentes.
No es una réplica de los LP antiguos DummyLp: cambian exportador, orden,
semillas, factores dyádicos y gap.

Gurobi 13.0.2, cuatro hilos, semillas 307,331,353, 12 s y 12 unidades de trabajo,
gap 0.1\%, tolerancias 1e-6/1e-5, sin starts. Factorial Base/Structural por
escalamiento interno auto/off y presolve auto/off, más controles GM y Ruiz.
Reserva con parámetros normales. Todas las resoluciones fueron secuenciales.
126 brazos, 108 llamadas MILP reales al deduplicar y 42 LP diagnósticos.
Un LP racional posterior se informa aparte. Diagnóstico y verificación no se
incluyen en el tiempo de optimización; sus costos están registrados.

\textbf{Atribución.} Un selector común elige potencias de dos dentro de una
ventana [1e-9,1e9], con banda identidad [0.5,2]. Se comprueba exactamente cada
producto de coeficiente/RHS y se conservan cotas, dominios y objetivo.
Algebraic usa GM; Structural cambia a L2 en acoplamientos. En los seis modelos,
ambos candidatos de acoplamiento quedan en la banda identidad y el resto usa
GM: \textbf{Algebraic=Structural en todas las filas}. Generator aplica además
presupuestos físicos: tampoco cambia factores. No se fuerza actividad reduciendo
presupuestos. Ruiz conserva columnas y usa cuatro medias iteraciones.

\begin{center}\begin{tabular}{lrrrr}\toprule
Escenario & MW & m$^3$ & Admisibles & Modificados/incompatibles\\\midrule
Exigente & .01 & 100 & 100\% & 0/0\\
Referencia & .1 & 1000 & 100\% & 0/0\\
Permisivo & 1 & 10000 & 100\% & 0/0\\\bottomrule
\end{tabular}\end{center}
Las unidades algebraicas se convierten por los coeficientes físicos de cada
fila; son objetivos numéricos del estudio, no tolerancias de UTE. La garantía
residuo original $\leq10^{-6}/\alpha$ es condicional al residuo exportado.
Los máximos observados divididos por presupuesto, sobre todos los puntos
finales/políticas, fueron 9.699e-05, 9.699e-06, 9.699e-07. Se evaluaron racionalmente y no se confunden
con la garantía. La igualdad de factores impide identificar una ventaja
operacional de la información física en este banco; no demuestra inutilidad
universal de esa información.

\newpage
\textbf{Resultado con escalamiento interno y presolve normales.}
Cada celda indica corridas que declaran OPTIMAL al gap pedido sobre tres, y
mediana de trabajo incluyendo las censuradas. GM agrupa Algebraic y Structural.

\begin{center}\begin{tabular}{llrrr}\toprule
Caso & Etapa & Base & GM & Ruiz\\\midrule
001 & D & 0/3; 12 & 3/3; 7.805 & 0/3; 12.01 \\
014 & D & 3/3; 0.453 & 3/3; 0.3963 & 3/3; 0.4278 \\
042 & D & 3/3; 0.3291 & 3/3; 0.4943 & 3/3; 0.3913 \\
002 & R & 0/3; 12.04 & 0/3; 12.13 & 0/3; 12 \\
003 & R & 0/3; 12 & 2/3; 4.931 & 0/3; 12 \\
004 & R & 3/3; 0.2682 & 3/3; 0.3541 & 3/3; 0.2581 \\
\bottomrule\end{tabular}\end{center}
D: descubrimiento; R: reserva. Una mediana de trabajo menor con peor cota o
incumbente no es una aceleración. RESULTADOS.md y los CSV conservan todos los
estados, gaps, cotas, objetivos, tiempos y el factorial completo. No se informa
una razón de aceleración cuando uno no termina o la calidad primal difiere.

\textbf{Qué revela presolve.} En 6 de seis casos Base y GM tienen
exactamente la misma huella de filas normalizadas, dominios y objetivo tras
el diagnóstico presolve normal, aunque sus magnitudes y trabajo difieren.
La relajación LP normal coincide numéricamente. Ni esa huella ni la reducción
de rangos permiten afirmar una ventaja universal. No se ha localizado un corte,
pivote o factor interno que explique causalmente las diferencias.

\textbf{Comprobación posterior.} Un testigo racional factible de costo 0 refuta 4 cotas inferiores del factorial; 4 provenían de corridas con estado OPTIMAL. Todas tenían ScaleFlag=0. Se verificaron sin tolerancias filas, cotas e integridad del testigo. Esto certifica una contradicción de cotas, no el óptimo del MILP.
El intento racional usa un patrón entero del descubrimiento, todas las filas y
cotas originales y costo $\leq0$, con 30 s de SoPlex. Está fuera de los tiempos
del factorial y es una validación posterior, no evidencia prospectiva.

\textbf{Límites y decisión.} Se fijaron conclusiones y hashes del descubrimiento
antes de abrir resultados de reserva. No se ajustó un predictor a tres casos
ni se reajustó con los reservados. Es un ensayo acotado de una familia y un
solver: no prueba una ventaja nueva generator-aware, un mecanismo interno
específico o una predicción útil fuera de muestra. Por eso se archiva separado
y no se amplía el manuscrito. El diagnóstico compacto LP--DW, la frontera
precisión--trabajo y los límites conjuntos del escalamiento siguen pendientes.

\textbf{Reproducción y fuentes.} El paquete incluye entradas, MPS, datos
canónicos, factores, puntos, trazas, logs, protocolo y scripts. Un replay sin
solver reimporta los seis MPS y comprueba los residuos con aritmética racional.
La entrega OMS y los siete paquetes anteriores se verifican por SHA256.
La documentación de \href{https://docs.gurobi.com/projects/optimizer/en/current/reference/python/model.html#Model.presolve}{Gurobi, Model.presolve}
advierte que su modelo puede diferir del usado por optimize: no es la matriz
realmente escalada internamente. Véanse también
\href{https://docs.gurobi.com/projects/optimizer/en/current/reference/parameters.html#parameter:ScaleFlag}{ScaleFlag}
y el \href{https://soplex.zib.de/doc/html/EXACT.php}{modo racional de SoPlex}.
\end{document}
